%!TeX root=../emmatop.tex
\chapter[Chapter \thechapter]{}
\lettrine[lines=4,lraise=0.3]{T}{he} next morning brought Mr Frank Churchill again. He came with Mrs Weston, to whom and to Highbury he seemed to take very cordially. He had been sitting with her, it appeared, most companionably at home, till her usual hour of exercise; and on being desired to chuse their walk, immediately fixed on Highbury.—»He did not doubt there being very pleasant walks in every direction, but if left to him, he should always chuse the same. Highbury, that airy, cheerful, happy-looking Highbury, would be his constant attraction.«—Highbury, with Mrs Weston, stood for Hartfield; and she trusted to its bearing the same construction with him. They walked thither directly.

Emma had hardly expected them: for Mr Weston, who had called in for half a minute, in order to hear that his son was very handsome, knew nothing of their plans; and it was an agreeable surprize to her, therefore, to perceive them walking up to the house together, arm in arm. She was wanting to see him again, and especially to see him in company with Mrs Weston, upon his behaviour to whom her opinion of him was to depend. If he were deficient there, nothing should make amends for it. But on seeing them together, she became perfectly satisfied. It was not merely in fine words or hyperbolical compliment that he paid his duty; nothing could be more proper or pleasing than his whole manner to her—nothing could more agreeably denote his wish of considering her as a friend and securing her affection. And there was time enough for Emma to form a reasonable judgment, as their visit included all the rest of the morning. They were all three walking about together for an hour or two—first round the shrubberies of Hartfield, and afterwards in Highbury. He was delighted with every thing; admired Hartfield sufficiently for Mr Woodhouse's ear; and when their going farther was resolved on, confessed his wish to be made acquainted with the whole village, and found matter of commendation and interest much oftener than Emma could have supposed.

Some of the objects of his curiosity spoke very amiable feelings. He begged to be shewn the house which his father had lived in so long, and which had been the home of his father's father; and on recollecting that an old woman who had nursed him was still living, walked in quest of her cottage from one end of the street to the other; and though in some points of pursuit or observation there was no positive merit, they shewed, altogether, a good-will towards Highbury in general, which must be very like a merit to those he was with.

Emma watched and decided, that with such feelings as were now shewn, it could not be fairly supposed that he had been ever voluntarily absenting himself; that he had not been acting a part, or making a parade of insincere professions; and that Mr Knightley certainly had not done him justice.

\begin{figure}[tbph]
\centering
\includegraphics[width=.9\linewidth]{24stopt}
\caption{He stopt to look in}
\end{figure}

Their first pause was at the Crown Inn, an inconsiderable house, though the principal one of the sort, where a couple of pair of post-horses were kept, more for the convenience of the neighbourhood than from any run on the road; and his companions had not expected to be detained by any interest excited there; but in passing it they gave the history of the large room visibly added; it had been built many years ago for a ball-room, and while the neighbourhood had been in a particularly populous, dancing state, had been occasionally used as such;—but such brilliant days had long passed away, and now the highest purpose for which it was ever wanted was to accommodate a whist club established among the gentlemen and half-gentlemen of the place. He was immediately interested. Its character as a ball-room caught him; and instead of passing on, he stopt for several minutes at the two superior sashed windows which were open, to look in and contemplate its capabilities, and lament that its original purpose should have ceased. He saw no fault in the room, he would acknowledge none which they suggested. No, it was long enough, broad enough, handsome enough. It would hold the very number for comfort. They ought to have balls there at least every fortnight through the winter. Why had not Miss Woodhouse revived the former good old days of the room?—She who could do any thing in Highbury! The want of proper families in the place, and the conviction that none beyond the place and its immediate environs could be tempted to attend, were mentioned; but he was not satisfied. He could not be persuaded that so many good-looking houses as he saw around him, could not furnish numbers enough for such a meeting; and even when particulars were given and families described, he was still unwilling to admit that the inconvenience of such a mixture would be any thing, or that there would be the smallest difficulty in every body's returning into their proper place the next morning. He argued like a young man very much bent on dancing; and Emma was rather surprized to see the constitution of the Weston prevail so decidedly against the habits of the Churchills. He seemed to have all the life and spirit, cheerful feelings, and social inclinations of his father, and nothing of the pride or reserve of Enscombe. Of pride, indeed, there was, perhaps, scarcely enough; his indifference to a confusion of rank, bordered too much on inelegance of mind. He could be no judge, however, of the evil he was holding cheap. It was but an effusion of lively spirits.

At last he was persuaded to move on from the front of the Crown; and being now almost facing the house where the Bateses lodged, Emma recollected his intended visit the day before, and asked him if he had paid it.

»Yes, oh! yes«—he replied; »I was just going to mention it. A very successful visit:—I saw all the three ladies; and felt very much obliged to you for your preparatory hint. If the talking aunt had taken me quite by surprize, it must have been the death of me. As it was, I was only betrayed into paying a most unreasonable visit. Ten minutes would have been all that was necessary, perhaps all that was proper; and I had told my father I should certainly be at home before him—but there was no getting away, no pause; and, to my utter astonishment, I found, when he (finding me nowhere else) joined me there at last, that I had been actually sitting with them very nearly three-quarters of an hour. The good lady had not given me the possibility of escape before.«

»And how did you think Miss Fairfax looking?«

»Ill, very ill—that is, if a young lady can ever be allowed to look ill. But the expression is hardly admissible, Mrs Weston, is it? Ladies can never look ill. And, seriously, Miss Fairfax is naturally so pale, as almost always to give the appearance of ill health.—A most deplorable want of complexion.«

Emma would not agree to this, and began a warm defence of Miss Fairfax's complexion. »It was certainly never brilliant, but she would not allow it to have a sickly hue in general; and there was a softness and delicacy in her skin which gave peculiar elegance to the character of her face.« He listened with all due deference; acknowledged that he had heard many people say the same—but yet he must confess, that to him nothing could make amends for the want of the fine glow of health. Where features were indifferent, a fine complexion gave beauty to them all; and where they were good, the effect was—fortunately he need not attempt to describe what the effect was.

»Well,« said Emma, »there is no disputing about taste.—At least you admire her except her complexion.«

He shook his head and laughed.—»I cannot separate Miss Fairfax and her complexion.«

»Did you see her often at Weymouth? Were you often in the same society?«

At this moment they were approaching Ford's, and he hastily exclaimed, »Ha! this must be the very shop that every body attends every day of their lives, as my father informs me. He comes to Highbury himself, he says, six days out of the seven, and has always business at Ford's. If it be not inconvenient to you, pray let us go in, that I may prove myself to belong to the place, to be a true citizen of Highbury. I must buy something at Ford's. It will be taking out my freedom.—I dare say they sell gloves.«

»Oh! yes, gloves and every thing. I do admire your patriotism. You will be adored in Highbury. You were very popular before you came, because you were Mr Weston's son—but lay out half a guinea at Ford's, and your popularity will stand upon your own virtues.«

They went in; and while the sleek, well-tied parcels of »Men's Beavers« and »York Tan« were bringing down and displaying on the counter, he said—»But I beg your pardon, Miss Woodhouse, you were speaking to me, you were saying something at the very moment of this burst of my amor patriae. Do not let me lose it. I assure you the utmost stretch of public fame would not make me amends for the loss of any happiness in private life.«

»I merely asked, whether you had known much of Miss Fairfax and her party at Weymouth.«

»And now that I understand your question, I must pronounce it to be a very unfair one. It is always the lady's right to decide on the degree of acquaintance. Miss Fairfax must already have given her account.—I shall not commit myself by claiming more than she may chuse to allow.«

»Upon my word! you answer as discreetly as she could do herself. But her account of every thing leaves so much to be guessed, she is so very reserved, so very unwilling to give the least information about any body, that I really think you may say what you like of your acquaintance with her.«

»May I, indeed?—Then I will speak the truth, and nothing suits me so well. I met her frequently at Weymouth. I had known the Campbells a little in town; and at Weymouth we were very much in the same set. Colonel Campbell is a very agreeable man, and Mrs Campbell a friendly, warm-hearted woman. I like them all.«

»You know Miss Fairfax's situation in life, I conclude; what she is destined to be?«

»Yes—(rather hesitatingly)—I believe I do.«

»You get upon delicate subjects, Emma,« said Mrs Weston smiling; »remember that I am here.—Mr Frank Churchill hardly knows what to say when you speak of Miss Fairfax's situation in life. I will move a little farther off.«

»I certainly do forget to think of her,« said Emma, »as having ever been any thing but my friend and my dearest friend.«

He looked as if he fully understood and honoured such a sentiment.

When the gloves were bought, and they had quitted the shop again, »Did you ever hear the young lady we were speaking of, play?« said Frank Churchill.

»Ever hear her!« repeated Emma. »You forget how much she belongs to Highbury. I have heard her every year of our lives since we both began. She plays charmingly.«

»You think so, do you?—I wanted the opinion of some one who could really judge. She appeared to me to play well, that is, with considerable taste, but I know nothing of the matter myself.—I am excessively fond of music, but without the smallest skill or right of judging of any body's performance.—I have been used to hear her's admired; and I remember one proof of her being thought to play well:—a man, a very musical man, and in love with another woman—engaged to her—on the point of marriage—would yet never ask that other woman to sit down to the instrument, if the lady in question could sit down instead—never seemed to like to hear one if he could hear the other. That, I thought, in a man of known musical talent, was some proof.«

»Proof indeed!« said Emma, highly amused.—»Mr Dixon is very musical, is he? We shall know more about them all, in half an hour, from you, than Miss Fairfax would have vouchsafed in half a year.«

»Yes, Mr Dixon and Miss Campbell were the persons; and I thought it a very strong proof.«

»Certainly—very strong it was; to own the truth, a great deal stronger than, if I had been Miss Campbell, would have been at all agreeable to me. I could not excuse a man's having more music than love—more ear than eye—a more acute sensibility to fine sounds than to my feelings. How did Miss Campbell appear to like it?«

»It was her very particular friend, you know.«

»Poor comfort!« said Emma, laughing. »One would rather have a stranger preferred than one's very particular friend—with a stranger it might not recur again—but the misery of having a very particular friend always at hand, to do every thing better than one does oneself!—Poor Mrs Dixon! Well, I am glad she is gone to settle in Ireland.«

»You are right. It was not very flattering to Miss Campbell; but she really did not seem to feel it.«

»So much the better—or so much the worse:—I do not know which. But be it sweetness or be it stupidity in her—quickness of friendship, or dulness of feeling—there was one person, I think, who must have felt it: Miss Fairfax herself. She must have felt the improper and dangerous distinction.«

»As to that—I do not\longdash«

»Oh! do not imagine that I expect an account of Miss Fairfax's sensations from you, or from any body else. They are known to no human being, I guess, but herself. But if she continued to play whenever she was asked by Mr Dixon, one may guess what one chuses.«

»There appeared such a perfectly good understanding among them all\longdash« he began rather quickly, but checking himself, added, »however, it is impossible for me to say on what terms they really were—how it might all be behind the scenes. I can only say that there was smoothness outwardly. But you, who have known Miss Fairfax from a child, must be a better judge of her character, and of how she is likely to conduct herself in critical situations, than I can be.«

»I have known her from a child, undoubtedly; we have been children and women together; and it is natural to suppose that we should be intimate,—that we should have taken to each other whenever she visited her friends. But we never did. I hardly know how it has happened; a little, perhaps, from that wickedness on my side which was prone to take disgust towards a girl so idolized and so cried up as she always was, by her aunt and grandmother, and all their set. And then, her reserve—I never could attach myself to any one so completely reserved.«

»It is a most repulsive quality, indeed,« said he. »Oftentimes very convenient, no doubt, but never pleasing. There is safety in reserve, but no attraction. One cannot love a reserved person.«

»Not till the reserve ceases towards oneself; and then the attraction may be the greater. But I must be more in want of a friend, or an agreeable companion, than I have yet been, to take the trouble of conquering any body's reserve to procure one. Intimacy between Miss Fairfax and me is quite out of the question. I have no reason to think ill of her—not the least—except that such extreme and perpetual cautiousness of word and manner, such a dread of giving a distinct idea about any body, is apt to suggest suspicions of there being something to conceal.«

He perfectly agreed with her: and after walking together so long, and thinking so much alike, Emma felt herself so well acquainted with him, that she could hardly believe it to be only their second meeting. He was not exactly what she had expected; less of the man of the world in some of his notions, less of the spoiled child of fortune, therefore better than she had expected. His ideas seemed more moderate—his feelings warmer. She was particularly struck by his manner of considering Mr Elton's house, which, as well as the church, he would go and look at, and would not join them in finding much fault with. No, he could not believe it a bad house; not such a house as a man was to be pitied for having. If it were to be shared with the woman he loved, he could not think any man to be pitied for having that house. There must be ample room in it for every real comfort. The man must be a blockhead who wanted more.

Mrs Weston laughed, and said he did not know what he was talking about. Used only to a large house himself, and without ever thinking how many advantages and accommodations were attached to its size, he could be no judge of the privations inevitably belonging to a small one. But Emma, in her own mind, determined that he did know what he was talking about, and that he shewed a very amiable inclination to settle early in life, and to marry, from worthy motives. He might not be aware of the inroads on domestic peace to be occasioned by no housekeeper's room, or a bad butler's pantry, but no doubt he did perfectly feel that Enscombe could not make him happy, and that whenever he were attached, he would willingly give up much of wealth to be allowed an early establishment.